% ============================================================================
% Taller 1 — Datos espaciales, sistemas de referencia y georreferenciacion
% Estadistica Espacial (20929) · Corte I · Universidad El Bosque · 2026-II
%
% PLANTILLA DE ENTREGA. Se entrega el PDF compilado, por Brightspace.
%
% Como usarla:
%   1. Rellena la ficha de la portada. El digito de verificacion es lo primero
%      que hay que comprobar: un intervalo perfectamente calculado sobre las
%      estaciones equivocadas se ve identico a uno correcto.
%   2. Rellena la hoja de cifras de la pagina 2 ANTES de escribir prosa.
%   3. Responde cada literal donde dice \marcador{...}. Los marcadores en gris
%      que queden sin sustituir se ven en el PDF: esa es su unica funcion.
%   4. Compila con pdflatex (o subela a Overleaf tal cual: no usa ningun
%      paquete raro ni ningun archivo externo) y entrega el PDF.
%
% Nombre del archivo:  T1_Apellido_XXX.pdf   (XXX = tus tres ultimos digitos)
% ============================================================================
\documentclass[11pt,a4paper,oneside]{article}

\usepackage[utf8]{inputenc}
\usepackage[T1]{fontenc}
\usepackage[spanish,es-nodecimaldot]{babel}
\usepackage{lmodern}
\usepackage{microtype}
\usepackage[a4paper,top=2.3cm,bottom=2.3cm,left=2.4cm,right=2.4cm]{geometry}

\usepackage{graphicx}
\usepackage{booktabs}
\usepackage{array}
\usepackage{tabularx}
\usepackage{caption}
\usepackage{enumitem}
\usepackage{fancyhdr}
\setlength{\headheight}{14pt}
\usepackage{titlesec}
\usepackage{amsmath,amssymb}
\usepackage{xcolor}
\usepackage{listings}
\usepackage{float}
\usepackage{hyperref}

% --- Colores institucionales (los del informe del laboratorio) ---
\definecolor{ubosque}{RGB}{0,102,51}
\definecolor{grisclaro}{RGB}{240,240,240}
\definecolor{grismarcador}{RGB}{130,130,130}

\hypersetup{colorlinks=true, linkcolor=ubosque, urlcolor=blue!70!black,
            pdftitle={Taller 1 - Estadistica Espacial 20929}}

\pagestyle{fancy}
\fancyhf{}
\fancyhead[L]{\small\textit{Taller 1 --- Estadística Espacial (20929)}}
\fancyhead[R]{\small\textit{\leftmark}}
\fancyfoot[C]{\thepage}
\fancyfoot[R]{\small Universidad El Bosque}
\renewcommand{\headrulewidth}{0.4pt}
\renewcommand{\footrulewidth}{0.3pt}

\setlength{\parskip}{0.5\baselineskip}
\setlength{\parindent}{0pt}
\emergencystretch=3em
\titleformat{\section}{\normalfont\large\bfseries\color{ubosque}}{\thesection.}{0.5em}{}
\titleformat{\subsection}{\normalfont\normalsize\bfseries}{}{0em}{}
\captionsetup{font=small,labelfont=bf,labelsep=period}

\lstset{basicstyle=\ttfamily\footnotesize, breaklines=true, frame=single,
        rulecolor=\color{grisclaro!60!black}, columns=flexible,
        showstringspaces=false, xleftmargin=0pt, framexleftmargin=3pt}

% --- Los tres comandos de la plantilla ---
% Un hueco por rellenar. Se ve en gris y entre corchetes: si queda uno sin
% sustituir, salta a la vista en el PDF entregado.
% Ojo con la forma del color: `\textcolor{}{}` (comando) y NO
% `{\color{} ...}` (declaracion). Dentro de una celda de parrafo —las
% columnas X de tabularx— la declaracion arranca el parrafo con un whatsit y
% hunde su primera linea: la etiqueta de la fila quedaba arriba y su casilla
% media linea mas abajo, en la tabla de T4 y en la de la declaracion de IA.
% Aislado en un caso minimo: `{\itshape[x]}` alinea, `{\color{gris}[x]}` no.
\newcommand{\marcador}[1]{\textcolor{grismarcador}{\textit{[#1]}}}

\newcommand{\tarea}[4]{%
  \section[#1 · #2]{#1 · #2 \normalfont\small\color{grismarcador}(#3 \% del escrito · #4 palabras)}%
  \markboth{#1 · #2}{}}

% Un literal: la letra en negrita y el rotulo corto. El enunciado COMPLETO
% esta en el taller; aqui solo va el rotulo, para no gastar dos paginas en
% copiar lo que el estudiante ya tiene delante.
\newcommand{\literal}[2]{\vspace{0.35\baselineskip}\noindent\textbf{(#1)}~\textit{#2}\par}

\begin{document}

% ============================================================================
% PORTADA — la ficha de variante
% ============================================================================
\thispagestyle{empty}
\begin{center}

{\large\textbf{UNIVERSIDAD EL BOSQUE}}\\[0.15cm]
{\normalsize Facultad de Ciencias --- Estadística}\\[0.1cm]
{\normalsize Estadística Espacial (20929) --- 2026-II --- Corte I}\\[0.9cm]

\rule{\textwidth}{1.2pt}\\[0.25cm]
{\LARGE\textbf{Taller 1}}\\[0.15cm]
{\large\textit{Datos espaciales, sistemas de referencia\\ y georreferenciación}}\\[0.25cm]
\rule{\textwidth}{1.2pt}\\[0.9cm]

{\large\textbf{Trabajo individual}}\\[0.15cm]
{\normalsize Escrito 60\,\% \quad · \quad Defensa oral 40\,\%}\\[1.0cm]

\end{center}

\begin{tabular}{rl}
\textbf{Nombre completo:}          & \marcador{tu nombre} \\[0.25cm]
\textbf{Número de documento:}      & \marcador{completo} \\[0.25cm]
\textbf{Tus tres últimos dígitos:} & \marcador{000 a 999} \quad $\rightarrow$ \quad
                                     \textbf{variante} \marcador{la misma cifra} \\[0.25cm]
\textbf{Tu municipio:}             & \marcador{nombre} (\marcador{departamento}) \\[0.25cm]
\textbf{Llave en la MGN:}          & \texttt{\marcador{la que da la tira del taller}} \\[0.25cm]
\textbf{Tu patrón puntual:}        & número \marcador{NN} \quad con $n =$ \marcador{n} puntos \\[0.25cm]
\textbf{Dígito de verificación:}   & \marcador{suma de altitudes} m \\[0.25cm]
\textbf{Fecha de entrega:}         & \textbf{martes 25 de agosto de 2026, 9:00 a.\,m.} \\
\end{tabular}

\vspace{0.9cm}

\noindent\fcolorbox{ubosque}{grisclaro}{%
\begin{minipage}{0.95\textwidth}
\small
\textbf{Antes de escribir nada, comprueba el dígito de verificación.} Es la suma de las altitudes
de tus 40 estaciones. Si tu selección no da exactamente esa cifra, las 40 estaciones no son las
tuyas, y un intervalo perfectamente calculado sobre las estaciones equivocadas se ve idéntico a
uno correcto. Todo lo que construyas encima estaría mal sin que nada avise.
\end{minipage}}

\vspace{0.5cm}

\noindent\fcolorbox{grisclaro!60!black}{white}{%
\begin{minipage}{0.95\textwidth}
\small
\textbf{Cómo se entrega.} Un solo PDF por \textbf{Brightspace}, compilado con \LaTeX{} a partir de
esta plantilla. Nombre del archivo: \texttt{T1\_Apellido\_XXX.pdf}, donde \texttt{XXX} son tus tres
últimos dígitos. Máximo \textbf{13 páginas} sin contar el anexo. Las tablas y figuras no cuentan
para los límites de palabras de cada tarea.

\vspace{0.2cm}
\textbf{No se aceptan entregas tarde.} Lo que no esté en Brightspace a las 9:00 de ese martes no
se califica.
\end{minipage}}

\vfill
{\footnotesize\color{grismarcador}\today}

\newpage

% ============================================================================
% HOJA DE CIFRAS
% ============================================================================
\section{Hoja de cifras}
\markboth{Hoja de cifras}{}

Todas las cifras que el taller pide, juntas. \textbf{Rellénala antes de escribir prosa}: es lo que
permite que la discusión de cada tarea hable de números que ya están decididos, y es lo primero que
se revisa al calificar. Usa punto decimal y cinco decimales donde la cifra los tenga.

\vspace{0.3cm}
\renewcommand{\arraystretch}{1.25}
\begingroup\setlength{\parskip}{0pt}
\begin{tabularx}{\textwidth}{@{}l X r@{}}
\toprule
\textbf{Tarea} & \textbf{Cifra} & \textbf{Valor} \\
\midrule
T1 & Número de puntos $n$                                   & \marcador{~} \\
T1 & Área de la ventana                                     & \marcador{~} \\
T1 & Intensidad $\lambda$                                   & \marcador{~} \\
T1 & Distancia media observada al vecino más próximo $\bar d_{\min}$ & \marcador{~} \\
T1 & La que daría el azar $E[\bar d_{\min}]$                & \marcador{~} \\
T1 & Índice de Clark-Evans $R$                              & \marcador{~} \\
T1 & $R$ corregido por borde (Donnelly)                     & \marcador{~} \\
\midrule
T2 & Suma de altitudes (dígito de verificación), en m       & \marcador{~} \\
T2 & Media de temperatura, en $^\circ$C                     & \marcador{~} \\
T2 & Desviación estándar                                    & \marcador{~} \\
T2 & Error estándar clásico                                 & \marcador{~} \\
T2 & IC 95\,\% clásico                                      & \marcador{~} \\
T2 & $\rho$ (I de Moran con el vecino más próximo)          & \marcador{~} \\
T2 & $n_{\text{eff}}$                                       & \marcador{~} \\
T2 & Error estándar corregido                               & \marcador{~} \\
T2 & IC 95\,\% corregido                                    & \marcador{~} \\
\midrule
T5 & Distancia real entre tus dos primeras estaciones, en km & \marcador{~} \\
T5 & La misma, ya mal declarada, en «metros»                & \marcador{~} \\
T5 & Factor entre las dos                                   & \marcador{~} \\
T5 & Estaciones sin municipio en el \emph{join} mal declarado & \marcador{~} \\
T5 & Estaciones sin municipio en el \emph{join} correcto    & \marcador{~} \\
\bottomrule
\end{tabularx}
\endgroup

\vspace{0.45cm}
\noindent\textbf{T4 · los cuatro sistemas.} Es la tabla que pide el literal (a): las cuatro áreas,
las cuatro distancias, y cuánto mide de verdad un \emph{buffer} de 500 m construido en cada uno.

\vspace{0.25cm}
\begingroup\setlength{\parskip}{0pt}
\begin{tabularx}{\textwidth}{@{}l X X X@{}}
\toprule
\textbf{EPSG} & \textbf{Área (km$^2$)} & \textbf{Ancho de la caja (m)} & \textbf{\emph{Buffer} de 500 m: cuánto mide} \\
\midrule
4326 (s2) & \marcador{~} & \marcador{~} & \marcador{~} \\
3857      & \marcador{~} & \marcador{~} & \marcador{~} \\
3116      & \marcador{~} & \marcador{~} & \marcador{~} \\
9377      & \marcador{~} & \marcador{~} & \marcador{~} \\
\bottomrule
\end{tabularx}
\endgroup

\vspace{0.2cm}
{\small\color{grismarcador} T3, T6 y T7 no llevan cifras propias: sus datos vienen dados en el
enunciado y son los mismos para todo el curso.}

\newpage

% ============================================================================
% LAS SIETE TAREAS
% ============================================================================
\tarea{T1}{El régimen que no se ve}{10}{500}
\literal{a}{En qué régimen está tu patrón, justificado con las dos distancias y no con el índice solo. Si tu conclusión no se sostiene mirando el mapa, dilo.}
\marcador{tu respuesta}

\literal{b}{Qué \emph{no} prueba tu $R$: contesta con la escala a la que mira el índice, no con significancia.}
\marcador{tu respuesta}

\literal{c}{Un compañero tiene un $R$ casi igual con un $n$ muy distinto: ¿significan lo mismo? Justifica con la fórmula del denominador.}
\marcador{tu respuesta}

\literal{d}{Por qué la corrección de borde va en esa dirección y no en la otra, y si el cambio te movería la conclusión de (a).}
\marcador{tu respuesta}

\tarea{T2}{El intervalo que miente}{15}{600}
\literal{a}{La media, el IC 95\,\% clásico, $\rho$, $n_{\text{eff}}$ y el IC corregido, en una tabla.}
\marcador{remite a la hoja de cifras, y comenta aquí lo que la tabla no dice}

\literal{b}{Una sola frase, sin jerga, que un funcionario de la Secretaría de Ambiente pueda leer y usar.}
\marcador{una frase}

\literal{c}{¿El problema está en la estimación o en la confianza? Decide y justifica; «en las dos» no cuenta.}
\marcador{tu respuesta}

\literal{d}{Por qué la equicorrelación es falsa en el espacio y en qué dirección te engaña el resultado que calculaste.}
\marcador{tu respuesta}

\tarea{T3}{La auditoría}{20}{750}
\literal{a}{Cuál de los dos resultados está mal, con evidencia \textbf{interna a la tabla}: comprobable con los números que ves, sin recalcular nada.}
\marcador{tu respuesta, y la aritmética que la sostiene}

\literal{b}{Qué mide de verdad la columna mala, si no es lo que dice su encabezado.}
\marcador{tu respuesta}

\literal{c}{Cuantifica el daño en la banda donde más se note, y explica por qué justo ahí.}
\marcador{tu respuesta}

\literal{d}{Por qué el $I$ negativo del resultado correcto es lo esperable y no el error.}
\marcador{tu respuesta}

\literal{e}{Por qué este error no se vería mirando un mapa de las estaciones.}
\marcador{tu respuesta}

\tarea{T4}{Cuatro sistemas, una decisión}{15}{600}
\literal{a}{Las cuatro áreas y las cuatro distancias, en una tabla, diciendo cuál tomas como referencia y por qué.}
\marcador{tabla propia o remisión a la hoja de cifras, más la referencia elegida}

\literal{b}{Cuál eliges para trabajar, justificado \textbf{con la cifra de distorsión que mediste}. «Es el oficial de Colombia» no puntúa.}
\marcador{tu respuesta}

\literal{c}{En cuál de los cuatro un \emph{buffer} de 500 m mide de verdad 500 m, comprobado, y por qué en los otros no.}
\marcador{tu respuesta}

\literal{d}{Si la respuesta de (b) de un compañero con municipio en otra parte del país tiene que ser la misma, nombrando el origen de EPSG:3116.}
\marcador{tu respuesta}

\tarea{T5}{Reetiquetar no es reproyectar}{15}{600}
\literal{a}{Qué síntoma delata que algo va mal. Nómbralo antes de arreglarlo y sin usar la palabra «error»: di qué cifra concreta es imposible y por qué.}
\marcador{tu respuesta}

\literal{b}{Arréglalo: la misma medición antes y después, con qué función y por qué esa y no la otra.}
\marcador{tu respuesta}

\literal{c}{El \emph{join} espacial contra los municipios, mal declarado y bien, con cuántas estaciones se quedan sin municipio en cada caso.}
\marcador{tu respuesta}

\literal{d}{Si \texttt{st\_set\_crs} es tan peligroso, por qué existe: un caso real en que sí es la función correcta.}
\marcador{tu respuesta}

\tarea{T6}{La fuga}{10}{500}
\literal{a}{La diferencia entre las dos validaciones con el vocabulario de los capítulos: qué información tenía el modelo al azar que no tiene por bloques.}
\marcador{tu respuesta}

\literal{b}{Cuál de los dos $R^2$ reportarías a la Secretaría de Educación y qué pregunta responde cada uno. Las dos partes cuentan.}
\marcador{tu respuesta}

\literal{c}{Si al añadir una covariable esperarías que la brecha entre los dos $R^2$ creciera o se cerrara, y por qué.}
\marcador{tu respuesta}

\tarea{T7}{Auditar a la IA}{15}{600}
\literal{a}{Clasifica las tres afirmaciones como correcta, a medias o falsa, con una cifra o un contraejemplo detrás de cada veredicto.}
\marcador{tu respuesta}

\literal{b}{La falsa, reescrita para que quede correcta sin cambiar de tema.}
\marcador{tu reescritura}

\literal{c}{La de «a medias»: dónde se rompe exactamente, qué parte es cierta, cuál no, y qué cifra lo demuestra.}
\marcador{tu respuesta}

\literal{d}{Una frase: qué tipo de pregunta contesta bien un modelo sobre este material y cuál contesta mal, apoyada en lo que te pasó durante el taller.}
\marcador{una frase}

\newpage
% ============================================================================
% DECLARACIÓN DE USO DE IA
% ============================================================================
\section{Declaración de uso de inteligencia artificial}
\markboth{Declaración de uso de IA}{}

El uso de IA en este taller es \textbf{libre y no hay que pedir permiso}. Lo que se califica —el
10\,\% del escrito, dimensión de comunicación y honestidad— es esta declaración: qué consultaste y
\textbf{cómo lo verificaste}. Declarar no penaliza; no declarar, sí.

Escribe una fila por tarea. Si en alguna no consultaste nada, escríbelo: «no consulté» es una
declaración válida y dejarla en blanco no lo es.

\vspace{0.3cm}
\renewcommand{\arraystretch}{1.4}
\begingroup\setlength{\parskip}{0pt}
\begin{tabularx}{\textwidth}{@{}l X X@{}}
\toprule
\textbf{Tarea} & \textbf{Qué consulté y a qué modelo} & \textbf{Cómo verifiqué la respuesta} \\
\midrule
T1 & \marcador{~} & \marcador{~} \\
T2 & \marcador{~} & \marcador{~} \\
T3 & \marcador{~} & \marcador{~} \\
T4 & \marcador{~} & \marcador{~} \\
T5 & \marcador{~} & \marcador{~} \\
T6 & \marcador{~} & \marcador{~} \\
T7 & \marcador{~} & \marcador{~} \\
\bottomrule
\end{tabularx}
\endgroup

% ============================================================================
% DECISIONES PARA LA DEFENSA
% ============================================================================
\section{Decisiones que sostendré en la defensa}
\markboth{Decisiones para la defensa}{}

La defensa vale el 40\,\%, y hay una regla que le da dientes: \textbf{una decisión entregada por
escrito que no se pueda sostener en la defensa se recalifica}. Escribe entre tres y cinco
decisiones tuyas —no resultados: decisiones— y de qué tarea salen.

\begin{enumerate}[leftmargin=*,itemsep=4pt]
  \item \marcador{decisión} \hfill \textcolor{grismarcador}{\small (tarea \marcador{Tn})}
  \item \marcador{decisión} \hfill \textcolor{grismarcador}{\small (tarea \marcador{Tn})}
  \item \marcador{decisión} \hfill \textcolor{grismarcador}{\small (tarea \marcador{Tn})}
  \item \marcador{opcional} \hfill \textcolor{grismarcador}{\small (tarea \marcador{Tn})}
  \item \marcador{opcional} \hfill \textcolor{grismarcador}{\small (tarea \marcador{Tn})}
\end{enumerate}

\newpage
% ============================================================================
% ANEXO — no cuenta para el límite de páginas
% ============================================================================
\section*{Anexo. Código ejecutado y entorno}
\markboth{Anexo}{}
\addcontentsline{toc}{section}{Anexo}

Pega aquí el código que ejecutaste de verdad para T2, T4 y T5, y la salida de
\texttt{sessionInfo()}. No se califica su elegancia: sirve para saber sobre qué 40 estaciones
trabajaste si alguna cifra no cuadra.

\begin{lstlisting}[language=R]
# Pega aqui tu codigo
\end{lstlisting}

\begin{lstlisting}
# Pega aqui la salida de sessionInfo()
\end{lstlisting}

\end{document}
